% section_2_framework.tex — Conceptual Framework + Empirical Strategy
% Drafted per: docs/Draft/THESIS_SKELETON.md v5.1 §II (post-capex-drop 2026-04-22):
%   §2.1 Pre-commitment statement
%   §2.2 Conceptual framework (uncertainty and corporate decisions)
%   §2.3 Speech uncertainty measurement (DWZ + BGT + CEO partition + CEO-as-face)
%   §2.4 Precautionary motive + financing conservatism (HC + HFC anchor)
%   §2.5 Empirical design
% Pending task: T44 (Draft §II Framework + Empirical Strategy)

\section{Conceptual Framework and Empirical Strategy}
\label{sec:framework}

\subsection{Pre-Commitment Statement}
\label{sec:framework:precommit}

This section front-loads the statistical conventions used throughout the paper so that the empirical results in \S\ref{sec:main} and \S\ref{sec:additional} can be read with the inference rules held constant.

Two formal hypotheses are stated in \S\ref{sec:main}: HC (Cash---precautionary liquidity buffer, \S\ref{sec:main:hc}) and HFC (Financial Constraint moderation of HC, \S\ref{sec:main:hfc}). Both are tested one-tailed in the theoretically predicted positive direction on both CEO-specific speech-uncertainty regressors, UncAnsCEO and UncPreCEO. The one-tailed convention is justified by the sign restriction imposed by the precautionary motive of \citeA{opler1999} and by the dominant theoretical direction in the related precautionary-motive and financing-conservatism literature. One-tailed inference applies symmetrically to both regressors; the Q\&A and Presentation segments of CEO speech are both hypothesized to carry precautionary uncertainty content, and any channel-specific asymmetry observed empirically is reported descriptively rather than asserted ex ante.

Two additional analyses in \S\ref{sec:additional} are tested under the same one-tailed positive convention but are not formal hypotheses. The driver regressions of \S\ref{sec:additional:drivers} invert the direction of the main-suite specifications to ask whether the CEO speech-uncertainty metric itself responds to independently measured exogenous uncertainty shocks; we report them as construct-validation evidence for the main regressor. The outside-world reaction regressions of \S\ref{sec:additional:reaction} ask whether post-call bid-ask spreads widen with CEO speech uncertainty, consistent with the canonical information-asymmetry prediction that liquidity providers widen quoted spreads to compensate for perceived informed-trading risk. Neither set of regressions is used as an identification strategy for HC or HFC.

All specifications reported in the paper body appear in their corresponding tables in their entirety; no column is dropped between estimation and reporting. The four-step fixed-effects ladder applied within each suite (\S\ref{sec:framework:design}) is a within-regressor robustness check, not an independent replication---the twelve columns of the HC table, for example, are not twelve independent tests of HC. Multiple-testing adjustment, consistent with the conventions above, is handled at the per-hypothesis rather than per-cell level.

\subsection{Conceptual Framework}
\label{sec:framework:conceptual}

A large literature documents that aggregate uncertainty influences corporate decisions on financing and investment. The leading aggregate measures---the news-based US Economic Policy Uncertainty index of \citeA{baker2016} and its global extension by \citeA{davis2016}, together with implied-volatility series such as the Cboe Volatility Index (VIX)---are sampled at monthly or daily frequency but vary uniformly across firms within a period. They identify the time-series of uncertainty but not its cross-sectional incidence at a given firm.

A complementary literature recovers firm-specific uncertainty content from the firm's own corporate disclosures. Annual 10-K filings have been a primary source for textual measurement at firm-year frequency, with the canonical word-counting tradition introduced by \citeA{loughran2011} for the negative and uncertainty word lists. \citeA{hassan2020} extend the approach to earnings-call transcripts to construct a firm-quarter political-risk index, demonstrating that call speech contains firm-quarter-specific uncertainty content distinguishable from aggregate political shocks. \citeA{dzielinski2021} apply the \citeA{loughran2011} uncertainty word list to the same earnings-call corpus and show that uncertainty content varies across firms, across quarters within a firm, and across the prepared-remarks and Q\&A segments of a single call.

The natural unit of analysis for testing precautionary behavior at firm-quarter frequency is therefore the call-level speech-uncertainty measure: it is sampled quarterly, identifies the firm rather than the macroeconomy, and admits source-identifiable decomposition of speakers within a single call (CEO, CFO, other managers, analysts). The remainder of \S\ref{sec:framework} documents the speech-uncertainty construction we use (\S\ref{sec:framework:measurement}), the precautionary-motive theoretical anchor for the main hypotheses (\S\ref{sec:framework:precautionary}), and the empirical design (\S\ref{sec:framework:design}).
